\section*{Chapter \ref{ch:m1:asian-and-emerging-equity-markets} --- Asian and Emerging Equity Markets}

\begin{solution}{exo:m1:asian-and-emerging-equity-markets:1}
Upper limit 26.40, lower limit 21.60. The two agree here because
$24.00\times1.1$ falls on a tick. In general $P\times1.1$ does not, and
exchanges differ in rounding it to the nearest tick, inward or by a rule of
their own: a system that rounds differently sends orders one tick outside the
band, which are rejected exactly on the days that matter.
\end{solution}

\begin{solution}{exo:m1:asian-and-emerging-equity-markets:2}
$250\,000/61.30 = 4\,078$ shares, that is 8 full lots: 4\,000 shares, HKD
245\,200; duty $0.1\% = $ HKD 245.20 (and as much again on the eventual
sale). HKD 4\,800 stays uninvested: lot sizes matter for small accounts and
for precise hedges.
\end{solution}

\begin{solution}{exo:m1:asian-and-emerging-equity-markets:3}
(a) $0.1\%$ on each side: INR 5\,000 $+$ 5\,000 $=$ 10\,000. (b) $0.05\%
\times 20$ million $=$ INR 10\,000, on the sale only. (c) $0.15\% \times
400\,000 =$ INR 600.
\end{solution}

\begin{solution}{exo:m1:asian-and-emerging-equity-markets:4}
$\ln0.65/\ln0.9 = 4.09$: five days, four locked. At 20\%: $1.93$, two days,
one locked. At 5\%: $8.40$, nine days, eight locked.
\end{solution}

\begin{solution}{exo:m1:asian-and-emerging-equity-markets:5}
No tax: $h^\star = 6/4 = 1.5$ days. With duty: $26/4 = 6.5$ days. At ten
days the gross edge is 40 basis points: net 34 without the duty, 14 with it
--- the tax takes 59\% of the profit of a strategy it does not forbid.
\end{solution}

\begin{solution}{exo:m1:asian-and-emerging-equity-markets:6}
She cannot sell the shares. She can sell futures on the index, or on the
share if a single-stock future exists, for the same notional: that removes
the market part of the risk and leaves the share's specific move until
tomorrow's open, plus the basis of the hedge. If she already held the same
share from a previous day she could sell \emph{those} shares today: the ban
applies to the day's purchases, which is why local funds keep a standing
inventory in the names they trade intraday.
\end{solution}

\begin{solution}{exo:m1:asian-and-emerging-equity-markets:7}
Standard deviation and autocorrelation: 2.60\% and $+0.195$ at 5\%; 3.03\%
and $+0.088$ at 10\%; 3.26\% and $+0.033$ at 20\%, against a true 3.45\% and
zero. The tighter the limit, the more of the distribution is cut and carried
forward: volatility is understated and spurious momentum grows roughly in
proportion to the share of days that end at a limit.
\end{solution}

\begin{solution}{exo:m1:asian-and-emerging-equity-markets:8}
First, the close on a \omterm{def:m1:asian-and-emerging-equity-markets:limit}{limit-up} day is not a price at which one can buy: the
order book shows a queue of buyers and almost no sellers, and the expected
fill of a new order is close to zero --- and is largest precisely in the
cases where the next day is bad (when holders were willing to sell at the
limit). Second, the purchase could not be sold on the day it is made, so the
``next-day'' exit itself may be locked if the share gaps down. The statistic
to report is the return conditional on an order \emph{actually filled} at the
limit, weighted by filled quantity: a fill-weighted, not an equal-weighted,
average.
\end{solution}

\begin{solution}{pb:m1:asian-and-emerging-equity-markets:1}
\textbf{1.} $29.00/(1 + 0.04\times0.25) = 28.71$.
\textbf{2.} 22.00; 24.20; 26.62; 29.28.
\textbf{3.} Friday: 29.28 is the first limit above 28.71. Locked Tuesday,
Wednesday and Thursday: three days.
\textbf{4.} $\ln(1.4356)/\ln(1.1) = 3.79$; $\lceil3.79\rceil - 1 = 3$.
\textbf{5.} None: 9 million shares are ahead of you and 300\,000 trade.
\textbf{6.} $28.71/22.00 - 1 = 30.5\%$.
\textbf{7.} $100\,000 \times (28.71 - 22.00) = 671\,000$.
\textbf{8.} Time priority at the opening: the orders entered in the first
microseconds after order entry opens, or queued in the pre-open by the
brokers with the fastest lines. A firm would pay up to the expected gain
times its probability of being filled --- which is why \omterm{def:m1:asian-and-emerging-equity-markets:limit}{limit-up} queues are a
latency race in markets with \omterm{def:m1:asian-and-emerging-equity-markets:limit}{price limits}.
\textbf{9.} Wednesday at the earliest, at Wednesday's limit of 24.20 at best
--- and only if someone buys there, which everyone will.
\textbf{10.} Closes 20.00, 22.00, 24.20, 26.62, 28.71: returns $+10\%$,
$+10\%$, $+10\%$, $+7.9\%$.
\textbf{11.} $\sum r^2 = 0.0362$ against $(0.4356)^2 = 0.1898$: five times
less.
\textbf{12.} $\sqrt{0.0362/20} = 4.3\%$ a day, against $\sqrt{0.1898/20} =
9.7\%$ for the same move made in one day.
\textbf{13.} $(0.01 + 0.01 + 0.0079)/0.0362 = 0.77$: a momentum signal sees
the strongest trend in its sample, in a share nobody could buy.
\textbf{14.} $29.00/28.60 - 1 = 1.40\%$ in three months, 5.6\% a year.
\textbf{15.} Less 0.10\% of duty and 1.00\% of financing: 0.30\% in three
months, 1.2\% a year above the interest rate. With completion certain, that
is a pure liquidity premium; whether it is worth the balance sheet depends on
the firm's hurdle.
\textbf{16.} $2\,000\,000/28.60 = 69\,930$: 699 lots, 69\,900 shares.
\textbf{17.} It would have opened near 28.7 (in its own currency) on Tuesday.
During the three locked days the gap between the two lines measures the
move the limited market has not yet been allowed to make, not an arbitrage:
the lines cannot be delivered against each other.
\textbf{18.} Deal-break risk. With a probability $p$ of failure and a fall to
about 20 in that case, the 1.40\% gross return must pay for an expected loss
of $p \times 30\%$: at $p = 5\%$ the trade already loses money.
\textbf{19.} \textbf{Three days.}
\textbf{20.} It protects participants from trading at prices formed in panic
or error within a day, at the cost of suspending everyone's ability to trade
at the right price for as many days as the news is large.
\end{solution}

\begin{solution}{iq:m1:asian-and-emerging-equity-markets:1}
A band around the previous close outside which no order is accepted. If fair
value moves beyond it, the share goes \omterm{def:m1:asian-and-emerging-equity-markets:limit}{limit-up} or limit-down: one-sided queue,
almost no trading, and the move is completed over several days, each day's
band being reset from the previous close.
\iqlookfor{the one-sided queue, and that the move still happens.}
\end{solution}

\begin{solution}{iq:m1:asian-and-emerging-equity-markets:2}
Volatility (understated, since large moves are split across days) and
autocorrelation (spuriously positive). Also tail estimates and value at risk,
correlations with unconstrained markets on event days, and any backtest that
assumes trading at a limit-day close.
\iqlookfor{the mechanism, carry-over, not just the list.}
\end{solution}

\begin{solution}{iq:m1:asian-and-emerging-equity-markets:3}
Anything whose edge per round trip is below the 20 basis points of tax:
electronic market making, index and ETF arbitrage, intraday statistical
strategies. Liquidity is then provided by longer-horizon investors and by
brokers, spreads are wider, and price discovery migrates to instruments the
tax does not reach --- futures, swaps, depositary receipts or exempted \omterm{def:m1:what-a-trading-firm-does:market-maker}{market
makers}, where the law grants exemptions.
\iqlookfor{migration to untaxed instruments.}
\end{solution}

\begin{solution}{iq:m1:asian-and-emerging-equity-markets:4}
The two lines are not fungible: one cannot be delivered against the other,
and the investors allowed to hold each differ, as do their funding costs and
short-selling constraints. So there is no convergence date. It can be traded
as a statistical relative-value position --- long the cheap line, short the
dear one where shorting is possible --- sized for the possibility that the
gap widens for years, or as a view on access reforms that would narrow it.
\iqlookfor{no forced convergence, and sizing accordingly.}
\end{solution}

\begin{solution}{iq:m1:asian-and-emerging-equity-markets:5}
Rules as data, not code: a table keyed by market, instrument class and
validity date, holding lot, tick table, limit formula and its rounding,
session times, selling constraints; one generic validator that reads it.
Every change is a new dated row with its source circular; backtests query the
table as of the simulated date. Test with the exchange's own examples and
replay recent production orders against a new rule version before its
effective date. Rejections from the exchange are fed back as a monitoring
signal that the table is stale.
\iqlookfor{versioning by date and a single generic validator.}
\end{solution}

\begin{solution}{iq:m1:asian-and-emerging-equity-markets:6}
Most severe: long--short equity (no new shorts, existing shorts possibly
recalled), convertible-bond and warrant hedging, single-stock options market
making (cannot hedge sold puts or bought calls by shorting). Then index
arbitrage and ETF market making (the short-basket leg), and lending revenue
for long holders. Futures typically trade at a discount to fair value because
they become the only way to be short. First morning: check whether existing
shorts may be kept; replace single-name shorts by index futures where
correlation allows, accepting basis and specific risk; reduce \omterm{def:m1:pnl-and-positions:exposure}{gross exposure}
to what the hedge can support; and price in that the ban has no announced
end.
\iqlookfor{futures discount as a consequence, and concrete first actions.}
\end{solution}
